% Source: https://artofproblemsolving.com/community/c3h1663289_geometry_problem
% Stable identity: aops-topic-1663289-q01
Choose coordinates with $A_1$ at the origin and $A_1A_2$ horizontal. Each segment is perpendicular to the one before it, so the segments $A_1A_2,A_2A_3,\dots,A_{2014}A_{2015}$ are alternately horizontal and vertical. There are $2014$ segments: the $1007$ odd-numbered ones are horizontal and the $1007$ even-numbered ones are vertical, each of length $1$. Hence
\[
\overrightarrow{A_1A_{2015}}=\left(\varepsilon_1+\varepsilon_2+\dots+\varepsilon_{1007},\ \delta_1+\delta_2+\dots+\delta_{1007}\right),\qquad \varepsilon_i,\delta_i\in\{1,-1\},
\]
and every choice of the signs gives a valid configuration.

\textbf{Not determined.} If all signs are $+1$ (a staircase), $A_1A_{2015}=\sqrt{1007^2+1007^2}=1007\sqrt2$. Other choices give other values, as below.

\textbf{Least value.} A sum of $1007$ terms $\pm1$ is odd (it has the parity of $1007$), so it is not $0$ and its absolute value is at least $1$. Hence $A_1A_{2015}\ge\sqrt{1^2+1^2}=\sqrt2$. Equality occurs, for example, when the path always turns the same way: then every four steps go once around a unit square and return to the start, and since $2014=4\cdot503+2$ the path ends two steps (one horizontal, one vertical) from $A_1$. A path that does not retrace itself also works; the figure shows the same idea with $11$ points, where $5$ horizontal and $5$ vertical steps end at distance $\sqrt2$.

\begin{asy}
usepackage("amsmath");
size(4.5cm);
defaultpen(fontsize(10pt));
pair[] p={(0,0),(1,0),(1,1),(2,1),(2,2),(1,2),(1,3),(0,3),(0,2),(-1,2),(-1,1)};
for(int i=-1;i<=2;++i) draw((i,-0.3)--(i,3.3),gray(0.85));
for(int j=0;j<=3;++j) draw((-1.3,j)--(2.3,j),gray(0.85));
for(int i=0;i<10;++i) draw(p[i]--p[i+1],linewidth(1));
for(int i=0;i<11;++i) dot(p[i]);
draw(p[0]--p[10],red+dashed);
label("$A_1$",p[0],S);label("$A_2$",p[1],S);label("$A_{11}$",p[10],W);
label("$\sqrt2$",(p[0]+p[10])/2,SW,red);
label("\mbox{}",(min(currentpicture,user=true)+max(currentpicture,user=true))/2,black); // keeps the text after the figure black when LaTeX typesets the labels inline
\end{asy}

The least possible value is $A_1A_{2015}=\boxed{\sqrt2}$.
